A superlotação dos serviços de urgência é agravada pela presença de pacientes de baixo risco clínico, e a combinação de totens de autoatendimento e telemedicina pode oferecer a esses pacientes um canal alternativo dentro da própria unidade de saúde. Este trabalho propõe e avalia analiticamente uma arquitetura de software para telemedicina embarcada em totens, voltada ao atendimento pós-triagem de pacientes classificados como verde ou azul pelo Protocolo de Manchester em unidades públicas de saúde. A pesquisa é aplicada e segue o paradigma de Design Science Research, com revisão narrativa protocolada, derivação rastreável de requisitos, modelagem em camadas com o modelo C4 e avaliação analítica prévia por cenários de qualidade, verificação de conformidade normativa e análise de cobertura entre requisitos e componentes. O resultado é uma arquitetura em cinco camadas: totem, plataforma de teleconsulta, backend clínico, gateway HL7 FHIR R4 para a Rede Nacional de Dados em Saúde e o e-SUS Atenção Primária, e segurança transversal. A proposta inclui matriz de rastreabilidade com 19 requisitos, critérios clínicos de elegibilidade, reavaliação e retorno ao atendimento presencial, mapeamento das seções do Registro de Atendimento Clínico, especificação de autenticação, validação e tratamento de erros conforme a documentação oficial da rede, distinção entre base legal de tratamento de dados e consentimento assistencial, e controles de privacidade, contingência e acessibilidade. Na avaliação, todos os requisitos ficaram associados a componentes, 25 de 39 itens de conformidade foram atendidos em projeto, 9 parcialmente e 5 não puderam ser verificados sem protótipo ou homologação, e sete de nove cenários foram atendidos em projeto. Trata-se de proposta preliminar, não prototipada nem avaliada empiricamente, cuja compatibilidade é de projeto e depende de validação por especialistas, simulação do fluxo, homologação e análise de custo e desempenho.
